%! Choosing a Sample: Four Sampling Techniques — Unit 1, Lesson 1.3 — Warm-Up
% PILOT SHAPE (2026-09-06): modeled on AP Statistics lesson 1.4 minus its AP Practice page.
% ONE PAGE, blank and key. Three items at 12pt, \workrowsep 50pt, item spacing 22pt, so each
% item gets real handwriting room and still fits. Do not add a fourth item — the page is full.
% GRADUAL RELEASE: the warm-up rehearses, WITHOUT naming them, the three computations the four
% techniques run on. Item 1 is a part of a whole as a percent (spiral to 1.1). Item 2 is the
% stratified allocation with the name removed — section 2 of the notes opens on this 18. Item 3
% is the systematic sample with the name removed — k = 900/60 = 15, start 7, then 22 and 37 —
% which the second half of section 1 picks up. Same 900 students the notes use.
\documentclass[12pt]{article}
\usepackage{saar-article}
\usepackage{saar-boxes}
\newcommand{\TallMath}[1]{$\displaystyle #1\rule[-1.4em]{0pt}{3.2em}$}
% Word bank strip for every fill-in cluster (user request 2026-09-06). Defined per document;
% the key inherits it and prints the same strip, so heights match.
\newcommand{\wordbank}[1]{\par\vspace{2pt}\noindent\fcolorbox{goldacc}{goldbg}{\parbox{\dimexpr\linewidth-2\fboxsep-2\fboxrule\relax}{\small\textbf{Word bank:} #1}}\par\vspace{4pt}}
% Handwriting room inside every \begin{work} block. Moves the blank and the
% key together, so the two can never drift apart.
\setlength{\workrowsep}{50pt}

\begin{document}

\pageheader{Unit 1, Lesson 1.3}{Warm-Up}
% NAMESTRIP: no \namedateperiod here. The student writes their name once, on the
% cover this component is stapled behind.

\vspace{0.15in}

\begin{notesbox}{Spiral Review --- Riverbend High School}
\normalsize
Riverbend High School has $900$ students. The student council wants to survey
$60$ of them about a new late-bus schedule.

\begin{enumerate}[leftmargin=1.8em, itemsep=22pt, topsep=10pt]

  \item Of Riverbend's $900$ students, $270$ are in $9$th grade. What
        \textbf{percent} of the school is that?

        \begin{work}
          \dfrac{270}{900} &= 0.30 \\
                           &= 30\%
        \end{work}

  \item The council wants $9$th graders to make up the same percent of the
        \emph{survey} that they make up of the \emph{school}. How many of the
        $60$ students surveyed should be $9$th graders?

        \begin{work}
          0.30 \times 60 &= 18 \text{ students}
        \end{work}

  \item The office prints one list of all $900$ students, in order. You want to
        pick $60$ of them, spread evenly from the top of the list to the bottom.

        \textbf{(a)} How many names apart should the picks be?

        \begin{work}
          900 \div 60 &= 15
        \end{work}

        \textbf{(b)} You start at the $7$th name on the list and keep counting by
        your answer to (a). Which two names do you take next?
        \wordbank{8 \;\textbullet\; 9 \;\textbullet\; 15 \;\textbullet\; 22 \;\textbullet\; 37 \;\textbullet\; 52}

        \blank{2.2cm} \qquad \blank{2.2cm}

\end{enumerate}
\end{notesbox}

\end{document}
